\begin{table*}[t]
\centering
\caption{Original paper experiments. Positive changes favor the named profile over uniform. CE reduction is $100(L_U-L_F)/L_U$ using seed means; accuracy gain is in percentage points. Brackets give nominal paired 95\% confidence intervals (Fieller for CE; Student-$t$ for accuracy), conditional on the selected profiles and fixed data split. Profiles were selected by mean final test CE among the historical nonuniform runs, then frozen for the additional seeds; they stay fixed across columns. Pooled intervals do not account for the earlier selection. Final CE and accuracy use the last epoch. Minimum CE averages each seed's lowest recorded test loss; this retrospective endpoint uses test data to choose an epoch, not validation. The first two rows share the same uniform runs. See the main text for the assumptions and selection limits.}
\label{tab:loss_improvements}
\footnotesize
\setlength{\tabcolsep}{4pt}
\renewcommand{\arraystretch}{1.15}
\begin{tabular}{@{}lrlrrr@{}}
\toprule
Experiment & $n$ & Profile & \shortstack{Final CE\\reduction (\%)} & \shortstack{Min. test CE\\reduction (\%)} & \shortstack{Final accuracy\\gain (pp)} \\
\midrule
\shortstack[l]{CIFAR-10\\MLP schedules} & 25 & Step early & \shortstack{$+17.86$\\[-1pt]{\scriptsize $[17.31, 18.41]$}} & \shortstack{$+1.42$\\[-1pt]{\scriptsize $[1.19, 1.65]$}} & \shortstack{$+0.83$\\[-1pt]{\scriptsize $[0.55, 1.11]$}} \\
\shortstack[l]{CIFAR-10\\MLP budget controls} & 25 & Big step & \shortstack{$+22.61$\\[-1pt]{\scriptsize $[22.06, 23.15]$}} & \shortstack{$+2.07$\\[-1pt]{\scriptsize $[1.80, 2.33]$}} & \shortstack{$+1.08$\\[-1pt]{\scriptsize $[0.85, 1.31]$}} \\
\shortstack[l]{CIFAR-10\\ReLU p=0.1} & 10 & Big step & \shortstack{$+34.90$\\[-1pt]{\scriptsize $[34.36, 35.44]$}} & \shortstack{$+2.54$\\[-1pt]{\scriptsize $[2.03, 3.06]$}} & \shortstack{$+1.56$\\[-1pt]{\scriptsize $[1.22, 1.89]$}} \\
\shortstack[l]{CIFAR-10\\GELU p=0.1} & 10 & Big step & \shortstack{$+29.79$\\[-1pt]{\scriptsize $[28.96, 30.61]$}} & \shortstack{$+2.45$\\[-1pt]{\scriptsize $[1.90, 3.00]$}} & \shortstack{$+0.62$\\[-1pt]{\scriptsize $[0.16, 1.08]$}} \\
\shortstack[l]{CIFAR-100\\ViT} & 10 & Linear decreasing & \shortstack{$+4.15$\\[-1pt]{\scriptsize $[3.40, 4.90]$}} & \shortstack{$+1.39$\\[-1pt]{\scriptsize $[0.56, 2.21]$}} & \shortstack{$+0.66$\\[-1pt]{\scriptsize $[0.36, 0.96]$}} \\
\shortstack[l]{CIFAR-10\\ViT both-block ablation} & 10 & Step early & \shortstack{$+7.40$\\[-1pt]{\scriptsize $[5.35, 9.41]$}} & \shortstack{$+1.51$\\[-1pt]{\scriptsize $[0.33, 2.68]$}} & \shortstack{$+0.68$\\[-1pt]{\scriptsize $[0.24, 1.12]$}} \\
\bottomrule
\end{tabular}
\end{table*}


We ask two experimental questions: does moving a fixed amount of dropout toward the input improve performance, and does this advantage remain useful across different datasets and architectures? The CIFAR-10/100 experiments \cite{krizhevsky2009cifar} address the first question by holding the dropout budget fixed (App.~\ref{app:experimental_info}). The broader benchmarks address the second using each cohort's validation protocol for learning rate and dropout strength, allowing mean dropout to differ between profiles. The Vision Transformer experiments \cite{dosovitskiy2021vit} test empirical transfer; the MLP initialization theory does not establish their training dynamics. In the CIFAR comparisons, the largest gains appear at the end of training: final test loss falls by \(18\)--\(35\%\) in MLPs and \(4.15\)--\(7.40\%\) in Vision Transformers (Table~\ref{tab:loss_improvements}). Comparing each run's minimum recorded test loss gives smaller reductions of \(1.39\)--\(2.54\%\). This retrospective minimum uses the test set to choose an epoch; the broader benchmarks instead evaluate test performance at checkpoints chosen using validation loss.

Accuracy counts correct labels, while cross-entropy also measures how much probability they receive. In a binary example, reducing the probability assigned to the correct class from $0.8$ to $0.6$ leaves the prediction correct, but increases its loss from $0.22$ to $0.51$. Conversely, increasingly confident mistakes can outweigh improvements on other examples, allowing accuracy and loss to rise together. We therefore report both metrics. On CIFAR-100, the linearly decreasing profile improves final accuracy from $48.61\%$ to $49.27\%$, alongside the loss comparison.

\paragraph{Behavior after the loss minimum.}
In the 25-seed CIFAR-10 budget comparison (Fig.~\ref{fig:triple_dropout}), the mean test-loss curve under uniform dropout rises from $1.685$ at its minimum (epoch 15) to $2.223$ at epoch 75, whereas the early big-step profile rises from $1.650$ (epoch 29) to $1.721$. These increases are $32.0\%$ and $4.3\%$, respectively, and the big-step model's mean test accuracy improves from $41.88\%$ to $42.92\%$ over epochs 29--75. Front-loading therefore leaves room for classification accuracy to improve while limiting the late deterioration in cross-entropy. This pattern is consistent with reduced overfitting and better late-training generalization in this comparison; whether it comes from reduced memorization would require measuring memorization directly.

The appendices check that this is not a single-run effect: the gains persist across \(\bar h\)-sweeps and large-width sweeps for kinked \relu{} MLPs (App.~\ref{app:sweeps}, Fig.~\ref{fig:mlp_sweeps}), across the smooth GELU activation class (App.~\ref{app:sweeps}, Fig.~\ref{fig:gelu_h_sweep}), and across transformer schedules and component ablations (Figs.~\ref{fig:dropout_regimes_transformer} and~\ref{fig:ablations}). In these CIFAR experiments, the advantage weakens at high dropout or in narrow networks, consistent with the limitations of the small-field, large-width approximation.

\paragraph{Uncertainty across seeds.}
Tables~\ref{tab:loss_improvements} and~\ref{tab:frontloaded_results} report nominal $95\%$ confidence intervals, pairing runs by their recorded seed or original simulation index. Loss reductions use ratios of seed means, with paired Fieller intervals \cite{fieller1954interval}; accuracy gains use paired Student-$t$ intervals in percentage points. Both use $n-1$ degrees of freedom and assume independent seed pairs with approximately bivariate normal losses or normal accuracy differences. We extended the smaller headline comparisons with fresh seeds, freezing their schedules and hyperparameters beforehand: all original CIFAR comparisons now have at least ten pairs, as do all benchmark checkpoint comparisons. Benchmark final endpoints have five or ten pairs, as indicated in the table. These intervals describe seed variation on the recorded split, without adjustment for the earlier profile selection or multiple comparisons, and do not include uncertainty from new data or new splits. All six original CIFAR comparisons have positive minimum-loss and final-accuracy intervals; the MLP schedule and budget-control rows share their uniform runs and should not be counted as independent replications. The appendix retains the original multi-profile cohorts so that the broader comparisons remain visible.

\begin{table*}[t]
\centering
\caption{Frontloaded dropout across tasks. Profiles were selected by lowest mean minimum validation CE among complete paired arms in the historical cohorts, then frozen for the additional seeds. Each profile stays fixed across columns; pooled intervals do not account for the earlier selection. Positive values favor frontloading. CE reduction is $100(L_U-L_F)/L_U$ using seed means; accuracy gain is in percentage points. Brackets give nominal paired 95\% confidence intervals (Fieller for CE; Student-$t$ for accuracy); see the main text for assumptions and selection limits. Checkpoint metrics use each run's minimum-validation-loss checkpoint; final metrics use its last epoch. The $n$ column gives checkpoint/final (C/F) paired seed counts. Where historical final evaluations were absent, final results use only the five additional pairs. The extended Jannis Transformer uses weight decay $10^{-7}$ and fixed mean dropout $0.10$; other rows use zero weight decay and allow profile-specific dropout budgets. Learning rate and mean dropout follow each cohort's validation protocol. Per-epoch test histories were not recorded. Dashes mean missing final test evaluations, not zero change.}
\label{tab:frontloaded_results}
\footnotesize
\setlength{\tabcolsep}{4pt}
\renewcommand{\arraystretch}{1.15}
\begin{tabular}{@{}lrlrrrr@{}}
\toprule
& & & \multicolumn{2}{c}{Validation-selected test} & \multicolumn{2}{c}{Final-epoch test} \\
\cmidrule(lr){4-5}\cmidrule(l){6-7}
Experiment (training $N$) & \shortstack{$n$\\C/F} & Profile & \shortstack{CE reduction\\(\%)} & \shortstack{Accuracy gain\\(pp)} & \shortstack{CE reduction\\(\%)} & \shortstack{Accuracy gain\\(pp)} \\
\midrule
\shortstack[l]{Speech Commands (20,020)\\MLP} & 10/5 & Step early & \shortstack{$+12.56$\\[-1pt]{\scriptsize $[11.36, 13.73]$}} & \shortstack{$+1.72$\\[-1pt]{\scriptsize $[1.28, 2.16]$}} & \shortstack{$+12.59$\\[-1pt]{\scriptsize $[9.95, 15.13]$}} & \shortstack{$+1.71$\\[-1pt]{\scriptsize $[0.78, 2.64]$}} \\
\shortstack[l]{Speech Commands (20,020)\\Transformer} & 10/5 & Big step & \shortstack{$+9.04$\\[-1pt]{\scriptsize $[2.52, 15.24]$}} & \shortstack{$+2.06$\\[-1pt]{\scriptsize $[0.72, 3.41]$}} & \shortstack{$+15.51$\\[-1pt]{\scriptsize $[6.44, 23.04]$}} & \shortstack{$+1.10$\\[-1pt]{\scriptsize $[-0.10, 2.30]$}} \\
\addlinespace[3pt]
\shortstack[l]{Jannis (3,840)\\MLP} & 10/5 & Step early & \shortstack{$+2.60$\\[-1pt]{\scriptsize $[1.44, 3.74]$}} & \shortstack{$+1.56$\\[-1pt]{\scriptsize $[0.94, 2.19]$}} & \shortstack{$+16.23$\\[-1pt]{\scriptsize $[14.83, 17.60]$}} & \shortstack{$+3.03$\\[-1pt]{\scriptsize $[1.96, 4.11]$}} \\
\shortstack[l]{Jannis (3,840)\\Transformer} & 10/5 & Big step & \shortstack{$+0.15$\\[-1pt]{\scriptsize $[-0.62, 0.92]$}} & \shortstack{$+0.32$\\[-1pt]{\scriptsize $[-0.51, 1.15]$}} & \shortstack{$+13.02$\\[-1pt]{\scriptsize $[8.96, 16.81]$}} & \shortstack{$-0.19$\\[-1pt]{\scriptsize $[-1.36, 0.98]$}} \\
\shortstack[l]{Jannis extended (3,840)\\Transformer} & 10/-- & Step early & \shortstack{$-0.27$\\[-1pt]{\scriptsize $[-1.16, 0.61]$}} & \shortstack{$+0.39$\\[-1pt]{\scriptsize $[-0.31, 1.10]$}} & -- & -- \\
\addlinespace[3pt]
\shortstack[l]{Tiny ImageNet (20,000)\\MLP} & 10/5 & Step early & \shortstack{$+2.04$\\[-1pt]{\scriptsize $[1.72, 2.37]$}} & \shortstack{$+1.28$\\[-1pt]{\scriptsize $[0.99, 1.57]$}} & \shortstack{$+21.48$\\[-1pt]{\scriptsize $[19.59, 23.35]$}} & \shortstack{$+0.80$\\[-1pt]{\scriptsize $[0.47, 1.14]$}} \\
\shortstack[l]{Tiny ImageNet (20,000)\\Transformer} & 10/5 & Step early & \shortstack{$+2.21$\\[-1pt]{\scriptsize $[1.69, 2.73]$}} & \shortstack{$+1.49$\\[-1pt]{\scriptsize $[1.06, 1.92]$}} & \shortstack{$+21.06$\\[-1pt]{\scriptsize $[20.27, 21.83]$}} & \shortstack{$+0.73$\\[-1pt]{\scriptsize $[-0.19, 1.66]$}} \\
\shortstack[l]{Tiny ImageNet (80,000)\\MLP} & 10/10 & Step early & \shortstack{$+1.17$\\[-1pt]{\scriptsize $[0.92, 1.42]$}} & \shortstack{$+0.77$\\[-1pt]{\scriptsize $[0.50, 1.05]$}} & \shortstack{$+8.65$\\[-1pt]{\scriptsize $[8.26, 9.04]$}} & \shortstack{$+0.68$\\[-1pt]{\scriptsize $[0.48, 0.88]$}} \\
\shortstack[l]{Tiny ImageNet (80,000)\\Transformer} & 10/10 & Linear decreasing & \shortstack{$+0.76$\\[-1pt]{\scriptsize $[-0.24, 1.76]$}} & \shortstack{$+0.13$\\[-1pt]{\scriptsize $[-0.41, 0.67]$}} & \shortstack{$+1.66$\\[-1pt]{\scriptsize $[0.51, 2.81]$}} & \shortstack{$-0.00$\\[-1pt]{\scriptsize $[-0.41, 0.40]$}} \\
\addlinespace[3pt]
\shortstack[l]{FI-2010 (40,000)\\MLP} & 10/5 & Big step & \shortstack{$+4.25$\\[-1pt]{\scriptsize $[2.57, 5.88]$}} & \shortstack{$+0.81$\\[-1pt]{\scriptsize $[0.38, 1.25]$}} & \shortstack{$+68.23$\\[-1pt]{\scriptsize $[67.30, 69.12]$}} & \shortstack{$-2.48$\\[-1pt]{\scriptsize $[-2.86, -2.09]$}} \\
\bottomrule
\end{tabular}
\end{table*}


\subsection{Front-Loaded Dropout Across Datasets}
\label{sec:extended_benchmarks}

To test how far the scheduling advantage extends beyond CIFAR, we compare MLPs and Transformers on Speech Commands, Jannis, and Tiny ImageNet, together with a financial MLP on FI-2010. These comparisons span audio, tabular classification, and vision, using the depth-12, width-256 configurations detailed in App.~\ref{app:extended_experiments}. Table~\ref{tab:frontloaded_results} reports the test loss and accuracy of validation-selected models, alongside the final-epoch endpoints where they were recorded.

For each run, we retain the checkpoint with the lowest validation cross-entropy. We then choose the reported front-loaded profile by averaging these minimum validation losses across the paired seeds, restricting selection to profiles with complete seed coverage; test loss and accuracy enter only the subsequent comparison with uniform dropout. Most cohorts contain five seeds, while the extended Jannis Transformer comparison contains ten. This is a reanalysis of the saved validation histories, with the original hyperparameter searches and data splits held fixed. The selected profile stays the same across the checkpoint and final columns, so a change of endpoint cannot silently change the model being compared.

The clearest improvement occurs for the Speech Commands MLP: early-step dropout reduces checkpoint test loss by $12.71\%$ relative to uniform, with accuracy increasing from $72.23\%$ to $74.06\%$. Uniform dropout itself reduces loss by $10.82\%$ relative to no dropout, placing this comparison in a regime where regularization already helps. The Transformer gives a smaller, noisier improvement of $4.91\%$ with big-step dropout and an accuracy gain of $1.53$ percentage points; both confidence intervals include zero. For the MLP, the loss-reduction interval is $[11.01,14.36]\%$ and the accuracy-gain interval is $[1.20,2.46]$ percentage points. These audio results use a class-balanced random split of clips from the official training subset, leaving generalization to held-out speakers unresolved.

Jannis provides a tabular comparison with a different learning problem. Validation selects early-step dropout for the MLP, giving a $2.92\%$ reduction in checkpoint test loss and a $1.29$-percentage-point accuracy gain. For the standard Transformer cohort, validation selects big step, whose test-loss improvement is only $0.15\%$. The ten-seed follow-up is effectively a tie: its validation-selected early step increases test loss by $0.27\%$, despite improving accuracy by $0.39$ percentage points. Both Transformer loss intervals include zero, while the MLP interval is $[0.49,5.27]\%$, leaving the MLP as the clearer signal in this tabular comparison.

On Tiny ImageNet, the selected profiles reduce checkpoint loss by $1.94\%$ and $2.17\%$ for the MLP and ViT at 20,000 training examples, and by $1.26\%$ and $1.01\%$ at 80,000 examples. The larger-data MLP still reaches only $7.95\%$ accuracy, compared with $7.09\%$ for uniform, so its relative gain should be read alongside this weak absolute performance. At the final epoch, the corresponding loss reductions are $8.84\%$ and $1.68\%$; the difference from the checkpoint comparison makes clear how much of the final-loss gain develops after the validation optimum. For the 80,000-example ViT, both checkpoint and final loss intervals include zero. The FI-2010 MLP gives a $3.57\%$ checkpoint loss reduction with big-step dropout, with interval $[2.41,4.73]\%$; its $0.67$-percentage-point accuracy gain has an interval spanning zero.

These results compare complete training configurations whose learning rates and dropout strengths follow each cohort's validation protocol; mean dropout can therefore differ between the selected profile and uniform. The fixed-budget CIFAR experiments remain the direct control for allocation. The broader evidence shows a useful but task-dependent advantage, with small effects and losses retained alongside the positive cases. App.~\ref{app:extended_experiments} gives the learning curves and all recorded profiles, including the unrestricted nonuniform comparisons in Table~\ref{tab:benchmark_results}.

